% ===== 原文 pp.20–23（PDF p.30–33） =====
\chapter{Vafa-Witten 理论作为带规范场的 4d A-模型}

% [原文 p.20]
我们在上一章中已经断言：四维的 $\cN=4$ SYM 能以三种互不等价的方式扭转成 TFT，并在上一章末尾描述了其中一种，即 GL-扭转理论。本章中我们从一个新的视角来描述另一种扭转——Vafa-Witten 扭转理论。受三维带规范场 A-模型 TFT 的构造 [6] 的启发，我们在四维做类似的构造：一个带规范场的 A-模型 TFT；它可以看作一族 4d 拓扑规范理论（每选定一个靶空间 $X$，就有一个成员）。随后我们会把 Vafa-Witten 理论定位为这一族中的某个特定成员（取 $X=\fg$，即规范群的李代数）。

% [Tutor Note: p.20 原文此处特意括注 "not to be confused with the Donaldson-Witten twist of N=4"：本章的 Donaldson-Witten 理论指 Witten [45] 对 $\cN=2$ SYM 的扭转，而第二章提到的 Donaldson-Witten 扭转是 $\cN=4$ 的三种扭转之一（见 §2.5）。核对方法：$\cN=2$ 的扭转只用其 $U(1)_R$ 生成元与旋转群的组合，得到单一标量超荷 $Q$，满足 $\delta_Q^2=\delta_\sigma$（见 p.25 式 (3.2) 的说明）；而 $\cN=4$ 的扭转还有 GL-扭转参数 $t$ 这一整族。]

\section*{带规范场的拓扑 $\sigma$-模型}

上同调型 TFT 的最初例子是 Witten [45] 对 4d 中 $\cN=2$ SYM 的扭转版本（下文我们称之为 \emph{Donaldson-Witten} 理论，不要与上一章简要提及的 $\cN=4$ 的 Donaldson-Witten 扭转相混淆）。这是\emph{拓扑规范理论}的一个例子：其场含量包含一个规范场，它作为主 $G$-丛 $P$ 上的联络变换；此外还包含物质场，它们作为相伴向量丛的截面变换。

紧接 Donaldson-Witten 之后，又出现了 [10] 另一类上同调场论，即二维的 A-模型与 B-模型。它们是\emph{拓扑 $\sigma$-模型}的例子：玻色场 $\Phi$ 对应从 $M$ 到某个\emph{靶空间} $X$ 的映射，而物质场则作为拉回丛 $\Phi^*TX$ 的截面变换。

近年来，若干作者研究了一类混合型 TFT，它们弥合了上述两条发展脉络之间的空隙，因为它同时包含 $\sigma$-模型场与一个规范场。这些\emph{带规范场的拓扑 $\sigma$-模型}是这样构造的：取靶空间为一个允许李群 $G$ 作用的流形 $X$，并假设该作用是 $\sigma$-模型理论的整体对称性。然后通过把 $\sigma$-模型场与一个规范场耦合，把 $G$-作用提升为局域对称性；该规范场作为主 $G$-丛上的联络变换。最后，在作用量中再纳入一份拓扑规范理论，用以支配规范场的动力学。

% [Tutor Note: p.20 这里说 $G$ 在 $X$ 上的作用“被假设为 $\sigma$-模型理论的整体对称性”，这个条件并不自动成立：$\sigma$-模型的作用量含靶空间度量 $g_{IJ}$，故只有等距作用（$V_a$ 为 Killing 向量场）才是整体对称性。原文到 p.26 才明确写出“We take $l$ to act on $X$ by isometries”，前后呼应；核对方法：对 (3.4) 的 $\delta_\epsilon\phi^I=\epsilon^aV^I_a$ 检查 $\delta_\epsilon g_{IJ}=\epsilon^a(\nabla_IV_{aJ}+\nabla_JV_{aI})=0$。]

% [原文 p.21]
例如，[38] 中构造了三维拓扑 Rozansky-Witten $\sigma$-模型的带规范场版本，其规范扇区为 Chern-Simons 作用量。在 [4] 与 [15] 中，分别用 B-型与 BF 型规范扇区构造了 Rozansky-Witten 模型的其它规范化。在 [6] 中，构造了 A-模型的三维类似物，然后通过把它与一个 A-型规范理论耦合来完成规范化。

我们在四维构造一个新的带规范场拓扑 $\sigma$-模型。物质场是相伴纤维丛的截面，其典型纤维是一个弯曲流形 $X$。相伴纤维丛的数学或许不如相伴向量丛那样为人熟悉（尤其是作用量中的拓扑项与弯曲纤维上的等变上同调类之间的联系），因此我们在附录 C 中回顾这套数学。

% [Tutor Note: p.21 “作用量中的拓扑项与弯曲纤维上的等变上同调类之间的联系”正是附录 C 的主题。具体机制：把 $X$ 上的 $G$-等变闭形式拉回到全空间 $E$，再沿纤维积分，得到 $M$ 上的闭形式，其积分只依赖于上同调类，故不改变运动方程——这就是这类项称为“拓扑项”的原因。核对：与 §3.2.3 的 $W_{\cO}\in(S(\fg^*)\otimes\Omega^\bullet(X))^G$ 对照，$W_{\cO}$ 正是被拉回的对象。]

在 §3.1 中，我们首先回顾 4d A-模型，这是一个以任意实流形 $X$ 为靶空间的四维拓扑 $\sigma$-模型。接着在 §3.2 中，我们通过把 4d A-模型与一个规范群为 $G$ 的 Donaldson-Witten 规范理论耦合，来对它做规范化。最后在 §3.3 中我们证明：Vafa-Witten 扭转理论可以看作靶空间与 $G$-作用取特定选择时的带规范场 4d A-模型。这就把 Vafa-Witten 理论定位为拓扑规范理论模空间中的一个特定点。

\section{4d A-模型}

设 $X$ 是一个带黎曼度量 $g$ 的实流形，$M$ 是一个带黎曼度量 $h$ 的四维流形。我们构造一个底流形为 $M$、靶空间为 $X$ 的拓扑 $\sigma$-模型；即一个生活在 $M$ 上的拓扑场论，其场由映射 $\Phi:M\to X$ 以及若干在 $M$ 上取值于拉回丛 $\Phi^*TX$ 的形式组成。沿用 [6] 的术语，我们称该理论为以 $X$ 为靶空间的 \emph{4d A-模型}。“A-模型”这一命名源于它与熟悉的 2d A-模型（§7.1 回顾）之间的若干结构相似性。首先，它是一个上同调型拓扑 $\sigma$-模型。其次，与 2d A-模型一样（但与 2d B-模型不同），4d A-模型局域化在一组椭圆偏微分方程上。最后，在柱面 $S^1\times I$ 上约化时（其中 $S^1$ 是圆周，$I$ 是区间，两端施加特定的边界条件），4d A-模型约化为以 $T^*X$ 为靶空间的 2d A-模型。

在详细描述 4d A-模型之前，让我们稍微展开上一章中关于上同调型 TFT 的讨论。我们说过，上同调型 TFT 的作用量在费米型对称性 $\delta_Q$（BRST 变分）下不变，而 $\delta_Q$ 在模运动方程的意义下平方为零。

% [Tutor Note: p.21 作者给出 4d A-模型与 2d A-模型相似的三条理由，其中前两条可在本章内部核对：(i) 上同调型——有 $\delta_Q$ 且局域化；(ii) 局域化方程是椭圆方程——见 p.23 的 $\partial^\mu\phi^I+D_\nu B^{I\nu\mu}=0$ 与 p.24 关于主象征的论证（对照 §7.1 的 2d A-模型方程 $\dbar\phi=0$）。第三条（柱面 $S^1\times I$ 约化给出靶空间 $T^*X$）原文未给证明，需参考 [6]；阅读时可视其为待查断言，不宜自行补证。]

% [原文 p.22]
在我们将要讨论的所有模型中，我们都选择引入辅助场，使得即使不用运动方程也有 $\delta_Q^2=0$。此外，作用量具有 $U(1)$ \emph{鬼数}对称性；在该对称性下，$\delta_Q$ 的作用是对所有场增加一个单位的荷。

% [Tutor Note: p.22 “辅助场”的判据可对照 p.23 的作用量核验：$\tilde H$ 在 $S=\frac1{e^2}\delta_Q\int g\tilde\chi(\tilde H-2\partial\phi-2DB)$ 中只以代数方式出现（$\delta_Q\tilde H$ 中不含 $\tilde H$ 的导数），其运动方程 $0=\partial S/\partial\tilde H$ 是代数方程，故不带来新的动力学自由度——这正是“作用 $\delta_Q$ 后 $\tilde H$ 无动能项”的含义。]

% [Tutor Note: p.22 “$\delta_Q$ 的不动点轨迹”一语需要小心：$\delta_Q$ 是费米型（奇）算子，把玻色场变成费米场，故“被 $\delta_Q$ 固定”不可能按字面理解；实际含义是“$\delta_Q$ 变分为零的位形”，即原文紧接着说的“把 BRST 变分的右端对所有场置零”。用 (3.1) 核对：取 $\zeta=0$（于是 $\delta_Q\phi=0$）并要求 $\delta_QB^I=0$，即得 $\partial^\mu\phi^I+D_\nu B^{I\nu\mu}=0$ 这一局域化方程。]

一般而言，QFT 的关联函数是在无穷维场位形空间 $\cC$ 上的积分。TFT 的特殊之处在于：$\cC$ 上的积分局域化到一个有限维子空间 $\cM\subset\cC$，即费米型对称性 $\delta_Q$ 的不动点轨迹；把 BRST 变分 $\delta_Q$ 的右端对所有场置零，就得到一组\emph{局域化方程}。（对拓扑规范理论，空间 $\cM$ 还要用无穷维规范变换群对局域化子流形再取商。）因此，TFT 关联函数的研究就归结为模空间 $\cM$ 上各种积分的研究。

\subsection{场含量、作用量与变分}

4d A-模型 [6] 的玻色场包括一个映射
\[
  \Phi:M\to X
\]
以及场
\begin{gather*}
  B\in\Omega^{2+}(M,\Phi^*TX)\\
  \tilde H\in\Omega^1(M,\Phi^*TX),
\end{gather*}
其中记号 $B\in\Omega^{2+}(M,\Phi^*TX)$ 表示 $B$ 是 $M$ 上取值于拉回向量丛 $\Phi^*TX$ 的自对偶 2-形式。费米子包括
\begin{gather*}
  \zeta\in\Omega^0(M,\Phi^*TX)\\
  \tilde\chi\in\Omega^1(M,\Phi^*TX)\\
  \tilde\psi\in\Omega^{2+}(M,\Phi^*TX).
\end{gather*}
场 $\Phi$、$B$、$\tilde H$ 的鬼数为 0，而 $\zeta$ 与 $\tilde\psi$ 的鬼数为 1，$\tilde\chi$ 的鬼数为 $-1$。我们用局部函数 $\phi^I(x)$ 描述 $\Phi$ 的位形，其中 $\phi^I$ 是 $X$ 上的局部坐标，$I=1,\ldots,\dim X$，而 $x^\mu$ 是 $M$ 上的局部坐标，$\mu=1,\ldots,\dim M$。同样，取值于切丛的场按分量写作
\[
  \zeta=\zeta^I\partial_I
\]

% [原文 p.23]
其中 $\partial_I\equiv\partial/\partial\phi^I$，且 $\zeta^I=\zeta^I(x)$。

BRST 变分为
\begin{align}
  \delta_Q\phi^I &= \zeta^I \notag\\
  \delta_Q\zeta^I &= 0 \notag\\
  \delta_QB^I &= \tilde\psi^I-\Gamma^I_{JK}\zeta^JB^K \notag\\
  \delta_Q\tilde\psi^I &= \frac12R^I{}_{JKL}B^J\zeta^K\zeta^L-\Gamma^I_{JK}\zeta^J\tilde\psi^K \tag{3.1}\\
  \delta_Q\tilde\chi^I &= \tilde H^I-\Gamma^I_{JK}\zeta^J\tilde\chi^K \notag\\
  \delta_Q\tilde H^I &= \frac12R^I{}_{JKL}\tilde\chi^J\zeta^K\zeta^L-\Gamma^I_{JK}\zeta^J\tilde H^K \notag
\end{align}
其中 $\Gamma^I_{JK}$ 是 Christoffel 符号，$R^I{}_{JKL}$ 是关于度量 $g_{IJ}$ 的 Levi-Civita 联络的 Riemann 张量分量（都视为场 $\Phi$ 的函数）。对所有场成立的幂零条件 $\delta_Q^2=0$ 由 Riemann 张量的定义与 Bianchi 恒等式推出。

4d A-模型的作用量 $S$ 取为 BRST-正合形式
\[
  S=\frac{1}{e^2}\delta_Q\int_M \dd^4x\sqrt{h}\,g_{IJ}\tilde\chi^I_\mu\left(\tilde H^{J\mu}-2\partial^\mu\phi^J-2D_\nu B^{J\nu\mu}\right),
\]
其中
\[
  D_\nu B^{J\nu\mu}=\partial_\nu B^{J\nu\mu}+\gamma^\nu_{\nu\rho}B^{J\rho\mu}+\Gamma^J_{KL}\partial_\nu\phi^KB^{L\nu\mu}
\]
指的是 $\Lambda^2TM\otimes\Phi^*TX$ 上由 $h_{\mu\nu}$ 的 Levi-Civita 联络（系数 $\gamma^\mu_{\nu\rho}$）与 $g_{IJ}$ 的 Levi-Civita 联络的拉回共同诱导的联络。我们引入了一个乘在整个作用量前面的参数 $1/e^2$。由于它藏在 BRST-正合项内部，改变 $e$ 的值只会以 BRST-正合的量改变作用量，因此关联函数与 $e$ 的具体取值无关。作用 $\delta_Q$ 后可以发现：$\tilde H^I$ 没有动能项（即 $\tilde H^I$ 是辅助场），而其它玻色场的动能项取通常形式。

该理论的局域化方程由把上述 $\delta_Q$ 变分的右端置零给出，即方程
\[
  \partial^\mu\phi^I+D_\nu B^{I\nu\mu}=0
\]
（费米子置零）。这可以看作 2d A-模型的全纯瞬子方程（见 §7.1）的 4d 类比。对每个 $I$ 值，这个非线性偏微分方程代表四个方程、四个变量，因为自对偶性意味着对每个 $I$，$B^{I\nu\mu}$ 有三个独立分量，而 $\phi^I$ 对每个 $I$ 有一个分量。

% [Tutor Note: p.23 上式 $D_\nu B^{J\nu\mu}$ 中只出现一个 $\gamma$ 项，看似漏掉了 $\gamma^\mu_{\nu\rho}B^{J\nu\rho}$：该项因 $\gamma^\mu_{\nu\rho}$ 对 $\nu,\rho$ 对称（Levi-Civita 无挠）而 $B^{J\nu\rho}$ 反对称（2-形式），乘积恒为零，故原文写法正确。可用 $\gamma^\mu_{[\nu\rho]}=0$ 直接验证。]

% [Tutor Note: p.23 局域化方程的“四方程四变量”计数可逐 $I$ 核对：未知量为 $\phi^I$（1 个）与 $B^{I\nu\mu}$ 的 3 个独立分量（自对偶 $\Rightarrow$ $\Omega^{2+}$ 是 3 维），方程按 $\mu=0,1,2,3$ 恰为 4 个；这正是原文断言该方程组横向椭圆的前提。]
